\chapter{Ureteroscopy and Laser Lithotripsy}

\section*{Introduction \& Clinical Indications}
Ureteroscopy uses a rigid or flexible endoscope passed through the urethra and bladder into the ureter or kidney to remove or fragment urinary stones. Laser lithotripsy breaks stones into dust or removable fragments. The procedure is used for ureteric or renal stones that are unlikely to pass, cause obstruction or recurrent symptoms, or are unsuitable for observation or shock-wave lithotripsy. It may require more than one session for a large or complex stone burden.

An infected obstructed system is an emergency for drainage and antibiotics, not immediate definitive stone manipulation in an unstable patient. Ureteroscopy is selected according to stone size and location, anatomy, infection, renal function, pregnancy, anticoagulation, and patient preference. PCNL and shock-wave lithotripsy remain important alternatives.

\section*{Relevant Surgical Anatomy}
The ureter passes from the renal pelvis through the retroperitoneum to the bladder and has narrow points at the ureteropelvic junction, pelvic brim, and vesicoureteric junction. The ureteric orifice and intramural ureter must be entered without perforation. Flexible ureteroscopy follows the collecting system through the renal pelvis and calyces; access sheath placement is guided by ureteral caliber and risk of injury.

\section*{Preoperative Workup \& Preparation}
Assessment includes urinalysis and culture, pregnancy testing when relevant, renal function, blood count, coagulation and antithrombotic review, and noncontrast CT or suitable imaging. A positive culture is treated before elective stone manipulation. Consent covers ureteric injury, perforation, stricture, infection or sepsis, bleeding, residual fragments, stent symptoms, repeat procedures, and rare loss of renal function.

\section*{Surgical Technique \& Procedure}
Under anesthesia, a cystoscope identifies the ureteric orifice and a guidewire is placed. A semirigid or flexible ureteroscope is advanced under direct vision and fluoroscopy as needed. The stone is fragmented with laser energy, extracted with a basket when safe, and the collecting system is inspected. A ureteral stent may be left to maintain drainage after edema, trauma, infection, or extensive manipulation.

\section*{Complications \& Risks}
Risks include dysuria, hematuria, colic, urinary infection, sepsis, ureteral perforation or avulsion, false passage, retained fragments, ureteric stricture, stent migration or encrustation, and need for repeat surgery. Laser and basket use require careful protection of the ureter and avoidance of excessive intrarenal pressure.

\section*{Postoperative Care \& Recovery}
Pain, fever, urine output, hematuria, and stent symptoms are reviewed. The patient receives instructions about stent removal, hydration, analgesia, and urgent symptoms. Fever, rigors, severe flank pain, vomiting, reduced urine output, or inability to pass urine requires urgent assessment. Stone analysis and metabolic evaluation guide prevention.

\section*{Outcomes \& Prognosis}
Stone-free outcome depends on stone burden, location, anatomy, infection, access, and staged treatment. Ureteroscopy avoids a percutaneous tract but may require repeat procedures and stenting. Prevention of recurrence requires fluid, dietary, and medication advice based on stone composition and urine testing.

\section*{References}
\begin{enumerate}
  \item European Association of Urology. Guidelines on Urolithiasis.
  \item National Institute for Health and Care Excellence. Renal and Ureteric Stones: Assessment and Management.
  \item American Urological Association. Surgical Management of Stones: Guideline.
\end{enumerate}

\clearpage
